\documentclass[a4paper,12pt]{article}
\usepackage[T2A]{fontenc}
\usepackage[utf8]{inputenc}
\usepackage[russian]{babel}
\usepackage{amsmath,amssymb,mathtools}
\usepackage{helvet}
% Helvetica has no standard T2A shapes; use compatible Cyrillic sans-serif.
\DeclareFontFamily{T2A}{phv}{}
\DeclareFontShape{T2A}{phv}{m}{n}{<->ssub*cmss/m/n}{}
\DeclareFontShape{T2A}{phv}{b}{n}{<->ssub*cmss/bx/n}{}
\DeclareFontShape{T2A}{phv}{bx}{n}{<->ssub*cmss/bx/n}{}
\DeclareFontShape{T2A}{phv}{m}{it}{<->ssub*cmss/m/sl}{}
\renewcommand{\familydefault}{\sfdefault}
\usepackage{icomma}
\usepackage[margin=20mm,headheight=15pt,footskip=12mm]{geometry}
\usepackage[expansion=false]{microtype}
\usepackage{tikz}
\usetikzlibrary{arrows.meta,positioning,shapes.geometric,calc}
\usepackage{pgfplots}
\pgfplotsset{compat=1.18}
\usepackage{tabularx,booktabs}
\usepackage{enumitem,parskip,fancyhdr,setspace,xcolor}
\usepackage[unicode,hidelinks]{hyperref}
\definecolor{accent}{HTML}{527B79}
\definecolor{darkaccent}{HTML}{294D4B}
\definecolor{textgray}{HTML}{333333}
\color{textgray}
\onehalfspacing
\setlength{\parindent}{0pt}
\setlength{\parskip}{5pt}
\setlength{\emergencystretch}{2em}
\setlist[itemize]{leftmargin=7mm,itemsep=3pt,topsep=3pt}
\setlist[enumerate]{leftmargin=8mm,itemsep=7pt,topsep=4pt}
\renewcommand{\arraystretch}{1.25}
\pagestyle{fancy}
\fancyhf{}
\fancyhead[L]{\small Прикладные задачи с формулами}
\fancyhead[R]{\small Григорий Архипов}
\fancyfoot[C]{\small\thepage}
\renewcommand{\headrulewidth}{0pt}
\newcommand{\heading}[1]{\section*{\color{darkaccent}#1}}
\newcommand{\subheading}[1]{\subsection*{\color{darkaccent}#1}}
\newcommand{\checkitem}[1]{\par\noindent\makebox[6mm][l]{$\square$}\begin{minipage}[t]{\dimexpr\linewidth-6mm\relax}\raggedright #1\end{minipage}\par}
\tikzset{
 action/.style={rectangle,rounded corners=2mm,draw=accent,thick,align=center,text width=10cm,minimum height=13mm,font=\small,inner sep=3mm},
 terminal/.style={ellipse,draw=accent,thick,align=center,text width=6cm,minimum height=12mm,font=\small},
 decision/.style={diamond,aspect=3,draw=accent,thick,align=center,text width=5cm,font=\small,inner sep=1mm},
 side/.style={action,text width=3.4cm},
 arrow/.style={-{Stealth[length=3mm]},thick,draw=accent}
}
\hypersetup{pdftitle={Прикладные задачи с формулами — чек-лист},pdfauthor={Лёвин Артём Александрович}}
\begin{document}
\raggedright
\hypersetup{pageanchor=false}
\begin{titlepage}
\centering
\vspace*{22mm}
{\large Профильная математика\par}
\vspace{13mm}
{\Huge\bfseries Чек-лист\par}
\vspace{8mm}
{\LARGE Прикладные задачи\\с формулами\par}
\vspace{10mm}
{\large Решение уравнений\\и выбор подходящего корня\par}
\vspace{13mm}
{\large Григорий Архипов\par}
\vspace{6mm}
{\large 30 сентября 2026 года\par}
\vfill
{\normalsize Лёвин Артём Александрович\\эксклюзивно для Григория Архипова\par}
\vspace{14mm}
\begin{tikzpicture}
\draw[accent,line width=0.8pt] (0,0) .. controls (3,0.55) and (7,-0.55) .. (11,0);
\end{tikzpicture}
\vspace*{5mm}
\end{titlepage}
\hypersetup{pageanchor=true}

\heading{1. Как прочитать задачу и применить формулу}
Используйте чек-лист для повторения: отмечайте пункт, когда можете объяснить правило и применить его самостоятельно. Затем решите 10 задач Пути А. Дополнительно предложены 3 задачи на применение формул. Ответы приведены в конце пособия.

\checkitem{Мы различаем аргумент и значение функции. Запись $h(t)$ означает высоту в момент $t$. Скобки обозначают зависимость; умножение на $t$ записывают как $h\cdot t$.}
\checkitem{Мы определяем, что дано, что требуется найти и от какого уровня отсчитывается величина. Если стена имеет высоту $14$ м, а камень должен пройти на $1$ м выше неё, его высота над землёй равна $15$ м.}
\checkitem{Мы проверяем единицы, принятые в формуле, и смысловые ограничения: например, $t\ge0$, количество изделий целое, время ограничено продолжительностью процесса.}

\subheading{Единицы измерения и степени десяти}
\[
3\text{ мм}=0,003\text{ м}=3\cdot10^{-3}\text{ м},\qquad
1,2\cdot10^{-5}=12\cdot10^{-6}.
\]
При увеличении первого множителя в $10$ раз второй множитель должен уменьшиться в $10$ раз. Проверьте равенство численно. Десятичные дроби допустимы; выбирайте удобную форму записи.

\subheading{Пример с тепловым расширением}
Длина стержня при температуре $T$ задаётся формулой
\[
L(T)=L_0(1+\alpha T),\qquad
L_0=10\text{ м},\quad \alpha=1,2\cdot10^{-5}\,{}^{\circ}\text{C}^{-1}.
\]
Чтобы стержень удлинился на $3$ мм, должно выполняться
\[
L-L_0=\alpha L_0T=0,003,\qquad
T=\frac{0,003}{10\cdot1,2\cdot10^{-5}}=25\,{}^{\circ}\text{C}.
\]
Здесь $L_0$ соответствует $T=0\,{}^{\circ}\text{C}$. При иной исходной температуре в формуле расширения используют изменение температуры.

\clearpage
\heading{2. Порог и изменение величины}
\subheading{Находим границу через равенство}
\begin{center}
\begin{tabularx}{\linewidth}{@{}lX@{}}
\toprule
Фраза в условии & Математическая запись \\
\midrule
Не меньше; не менее $A$ & $f(x)\ge A$ \\
Не больше; не более $A$ & $f(x)\le A$ \\
Больше $A$; меньше $A$ & $f(x)>A$; $f(x)<A$ \\
Ровно $A$ & $f(x)=A$ \\
\bottomrule
\end{tabularx}
\end{center}
\checkitem{Мы находим граничное значение через уравнение $f(x)=A$, затем выбираем ответ по смыслу задачи. Такой способ подходит для рассмотренных задач на минимум, максимум и продолжительность.}

Фразы «не менее» и «не более» допускают равенство на границе. В этих задачах решаем уравнение для граничного значения, учитываем ход процесса и проверяем ответ. Для целого количества дополнительно выполняем округление в нужную сторону.

При цене $p$, переменных затратах $v$ на изделие и постоянных расходах $f$ прибыль равна $\Pi(q)=q(p-v)-f$.

\textbf{Пример.} При $p=600$ руб., $v=300$ руб. и $f=700\,000$ руб. прибыль от продажи $q$ изделий равна
$\Pi(q)=300q-700\,000$ руб. Для прибыли не менее $500\,000$ руб.:
\[
300q-700\,000=500\,000
\quad\Longleftrightarrow\quad q=4000.
\]
Прибыль растёт с числом проданных изделий, поэтому минимальный объём продаж равен $4000$ изделиям. Граничное значение допустимо: прибыль составляет ровно $500\,000$ руб. Если граница для целого количества равна $1666\frac23$, минимальное количество равно $1667$.

\subheading{Изменение значения функции}
\checkitem{Мы вычисляем значения в двух состояниях и вычитаем их в нужном порядке. Для нелинейной функции $f(t_1)-f(t_2)$ обычно отличается от $f(t_1-t_2)$.}

\textbf{Пример.} Расстояние падения камня до воды задаётся формулой $h(t)=5t^2$ в метрах. До дождя время равно $0,6$ с, после дождя --- $0,4$ с. Уровень воды поднялся на
\begin{align*}
\Delta h&=h(0,6)-h(0,4)=5(0,6^2-0,4^2)\\
&=5(0,6-0,4)(0,6+0,4)=1\text{ м}.
\end{align*}
Подстановка разности времён дала бы $5\cdot0,2^2=0,2$ м. Проверяйте, изменение какой величины требуется найти.

\clearpage
\heading{Алгоритм работы с прикладной формулой}
\thispagestyle{plain}
\begin{center}
\begin{tikzpicture}[node distance=9mm]
\node[terminal] (start) {Прочитайте условие};
\node[action,below=of start] (meaning) {Установите смысл обозначений,\ единицы и допустимые значения};
\node[action,below=of meaning] (model) {Переведите вопрос в формулу:\ найдите значение, изменение\ или составьте граничное уравнение};
\node[action,below=of model] (solve) {Выполните вычисления\ и найдите все кандидаты};
\node[action,below=of solve] (filter) {Проверьте кандидаты:\ область определения, ограничения формулы,\ слова «впервые», «наибольший», «не менее»};
\node[action,below=of filter] (check) {Подставьте результат в исходную формулу;\ проверьте единицы и смысл ответа};
\node[terminal,below=of check] (end) {Запишите требуемую величину};
\foreach \a/\b in {start/meaning,meaning/model,model/solve,solve/filter,filter/check,check/end}
\draw[arrow] (\a) -- (\b);
\end{tikzpicture}
\end{center}
\clearpage

\heading{3. Квадратное уравнение и отбор корней}
\checkitem{Мы избавляемся от числовых знаменателей умножением обеих частей на их общий кратный. Умножаем каждое слагаемое.}
\[
15=-\frac{x^2}{160}+\frac78x
\quad\Longleftrightarrow\quad
2400=-x^2+140x
\quad\Longleftrightarrow\quad
x^2-140x+2400=0.
\]
Наибольший знаменатель подходит, если остальные знаменатели являются его делителями. Например, для $6$ и $8$ можно взять $24$.

\checkitem{Мы решаем уравнение $ax^2+bx+c=0$, где $a\ne0$:}
\[
D=b^2-4ac,\qquad
x_{1,2}=\frac{-b\pm\sqrt D}{2a}\quad(D\ge0).
\]
При $D<0$ действительных корней нет; при $D=0$ есть один корень $-b/(2a)$. Если $b<0$, величина $-b$ положительна. Используйте скобки: $(-b)^2$, деление на $(-10)$.

\subheading{Пример с высотой над стеной}
Камень движется по закону $y=-x^2/160+7x/8$, где $x$ --- горизонтальное расстояние в метрах. На уровне $15$ м:
\[
D=140^2-4\cdot2400=10\,000,\qquad
x_1=20,\quad x_2=120.
\]
Оба значения дают $y=15$. Если требуется наибольшее расстояние, ответ равен $120$ м. Между этими точками камень летит выше $15$ м. Поэтому для наибольшего расстояния при условии «не менее $15$ м» также выбираем $120$ м: на границе равенство допустимо.

\checkitem{Мы сначала проверяем допустимость корней, затем выбираем ответ по вопросу задачи. «Впервые» задаёт самый ранний допустимый момент; «наибольшее» --- наибольшее допустимое значение.}

\textbf{Пример с нагревателем.} По формуле $T(t)=1400+200t-10t^2$ температура измеряется в кельвинах, время --- в минутах. При $T=1760$ получаем $t^2-20t+36=0$. Так как $\sqrt{256}=16$, корни равны $2$ и $18$. Прибор нужно отключить до опасного перегрева: предельное время работы --- $2$ мин.

\clearpage
\heading{4. Область применимости формулы}
\subheading{Доля объёма и доля высоты}
\checkitem{Мы заменяем долю объёма той же долей высоты для вертикального бака с постоянной площадью горизонтального сечения.}
\[
V=Sh,\quad V_0=Sh_0
\quad\Longrightarrow\quad
\frac{V}{V_0}=\frac{h}{h_0}.
\]
Для цилиндрического бака четверть начального объёма означает $h=h_0/4$. Для бака, площадь сечения которого меняется с высотой, нужна другая зависимость.

\subheading{Пример со сливом воды}
Высота воды в цилиндрическом баке при сливе задаётся формулой
\[
h(t)=20-\frac{t}{25}+\frac{t^2}{50\,000}
=20\left(1-\frac{t}{1000}\right)^2,
\qquad 0\le t\le1000,
\]
где $h$ измеряется в метрах, $t$ --- в секундах. При $t=1000$ бак пуст; после опустошения эта формула слива неприменима.

Для четверти начального объёма $h=5$ м. Поэтому
\begin{align*}
20-\frac{t}{25}+\frac{t^2}{50\,000}&=5,\\
t^2-2000t+750\,000&=0,\\
D&=1\,000\,000,\qquad \sqrt D=1000,\\
t_1&=500,\qquad t_2=1500.
\end{align*}
Корень $1500$ находится вне диапазона $[0;1000]$. Это решение уравнения недопустимо для описываемого слива. Ответ: $500$ с.

\checkitem{Мы проверяем границы применимости формулы до выбора корня. Продолжение параболы за момент опустошения бака создаёт математические значения, которые процесс слива уже не описывают.}

\clearpage
\heading{Алгоритм отбора корней}
\thispagestyle{plain}
\begin{center}
\begin{tikzpicture}[node distance=11mm]
\node[terminal] (start) {Найдите все корни уравнения};
\node[action,below=of start] (domain) {Оставьте корни из области определения\ и допустимого диапазона процесса};
\node[decision,below=of domain] (has) {Допустимые корни есть?};
\node[terminal,text width=3.4cm,right=10mm of has] (none) {Запишите:\ допустимых решений нет};
\node[action,below=14mm of has] (request) {Прочитайте, какую величину ищут:\ первый момент, наибольшее расстояние\ или продолжительность};
\node[action,below=of request] (select) {Выберите нужный корень\ или найдите разность границ\ $\Delta t=t_2-t_1$};
\node[terminal,below=of select] (end) {Проверьте и запишите ответ};
\draw[arrow] (start) -- (domain);
\draw[arrow] (domain) -- (has);
\draw[arrow] (has) -- node[above,font=\small]{нет} (none);
\draw[arrow] (has) -- node[right,font=\small]{да} (request);
\draw[arrow] (request) -- (select);
\draw[arrow] (select) -- (end);
\end{tikzpicture}
\end{center}
\clearpage

\heading{5. Граничные моменты и продолжительность}
\checkitem{Мы различаем моменты достижения уровня и продолжительность периода. Чтобы найти начало и конец периода, решаем уравнение для заданного уровня, затем вычитаем меньший момент из большего.}

\textbf{Пример.} Высота мяча задана формулой $h(t)=1,4+9t-5t^2$ в метрах, время $t$ --- в секундах после броска. Нужно определить, сколько времени высота составляет не менее $3$ м. Находим моменты достижения уровня $3$ м через равенство:
\begin{align*}
1,4+9t-5t^2&=3,\\
5t^2-9t+1,6&=0,\\
D&=81-32=49,\\
t_1=\frac{9-7}{10}&=0,2,\qquad t_2=\frac{9+7}{10}=1,6.
\end{align*}
Мяч достигает $3$ м при подъёме в момент $0,2$ с и при спуске в момент $1,6$ с. Между этими моментами он находится выше $3$ м. Поэтому требуемая продолжительность равна $1,6-0,2=1,4$ с. Условие «не менее» включает граничные моменты.

\begin{center}
\begin{tikzpicture}
\begin{axis}[
 width=12.5cm,height=5.7cm,
 axis lines=middle,xmin=0,xmax=2.1,ymin=0,ymax=6.2,
 xlabel={$t$, с},ylabel={$h$, м},
 xtick={0,0.2,0.9,1.6,2},ytick={0,1.4,3,5.45},
 xticklabels={0,{$0,2$},{$0,9$},{$1,6$},2},
 yticklabels={0,{$1,4$},3,{$5,45$}},
 samples=200,domain=0:1.94403,clip=false,
 tick label style={font=\small},label style={font=\small}]
\addplot[thick,color=accent]{1.4+9*x-5*x*x};
\addplot[dashed,color=textgray,domain=0:2]{3};
\addplot[only marks,mark=*,color=darkaccent] coordinates {(0.2,3) (1.6,3)};
\node[above right,font=\small] at (axis cs:0.2,3) {$t_1$};
\node[above left,font=\small] at (axis cs:1.6,3) {$t_2$};
\end{axis}
\end{tikzpicture}
\end{center}

\checkitem{Мы проверяем, что между найденными граничными моментами условие задачи выполняется. Для мяча это видно по графику: между двумя пересечениями с уровнем $3$ м парабола проходит выше этого уровня.}
\checkitem{Мы записываем ответ с единицами и отвечаем на вопрос: момент $t$, промежуток моментов или продолжительность $\Delta t$.}

\clearpage
\heading{6. Дополнительно: метод интервалов}
Граничные значения в рассмотренных задачах находим через уравнения. Метод интервалов помогает при необходимости уточнить подходящие промежутки и проверить включение границ.

\checkitem{Мы меняем знак неравенства при умножении или делении на отрицательное число. Например, $-t^2+8t-12\ge0$ равносильно $t^2-8t+12\le0$.}
\checkitem{Мы отмечаем нули и недопустимые точки многочлена или рационального выражения, затем определяем знак на каждом открытом промежутке. Пробное число выбираем внутри промежутка. Корень даёт ноль, поэтому для проверки знака не подходит.}

\textbf{Пример.} Для $p(t)=(t-2)(t-6)$:
\begin{center}
\begin{tabular}{@{}lccc@{}}
\toprule
Промежуток & $(-\infty;2)$ & $(2;6)$ & $(6;+\infty)$ \\
\midrule
Пробное число & $0$ & $4$ & $8$ \\
Знак $p(t)$ & $+$ & $-$ & $+$ \\
\bottomrule
\end{tabular}
\end{center}
В точках $2$ и $6$ выражение равно нулю. Следовательно,
\[
p(t)\le0\ \Longleftrightarrow\ t\in[2;6],\qquad
p(t)<0\ \Longleftrightarrow\ t\in(2;6).
\]
\checkitem{Мы включаем допустимые нули при знаках $\le$ и $\ge$. При строгих знаках $<$ и $>$ нули исключаем. Точки, где выражение не определено, исключаются всегда.}

\subheading{Скобки и типичные ошибки}
\begin{itemize}
\item $[a;b]$ включает обе границы; $(a;b)$ исключает обе. Смешанные скобки задают включение одной границы.
\item У бесконечности используют круглую скобку: $(-\infty;a]$.
\item В общем случае знаки проверяемт. Через корень чётной кратности знак не меняется; у двух различных простых корней квадратного трёхчлена знаки чередуются.
\item Если делите на выражение с переменной, сначала выясните его знак и допустимость деления. При делении на отрицательное число знак неравенства меняется.
\end{itemize}

\textbf{Пример с температурой.} Для $T(t)=-t^2+14t+25$ граничные моменты для уровня $70\,{}^{\circ}\text{C}$ находим из уравнения
\[
-t^2+14t+25=70
\ \Longleftrightarrow\ (t-5)(t-9)=0.
\]
Получаем $t_1=5$ и $t_2=9$ мин. Между ними температура выше $70\,{}^{\circ}\text{C}$, поэтому подходящие моменты: $t\in[5;9]$ мин. Продолжительность: $9-5=4$ мин.



\clearpage
\heading{Дополнительно: решение неравенства}
\thispagestyle{plain}
\begin{center}
\begin{tikzpicture}[node distance=9mm]
\node[terminal] (start) {Запишите ограничения};
\node[action,below=of start] (zero) {Перенесите всё в одну часть;\\ выберите запись $F(x)>0$, $F(x)\ge0$,\\ $F(x)<0$ или $F(x)\le0$};
\node[action,below=of zero] (roots) {Найдите нули и точки,\ где выражение не определено};
\node[action,below=of roots] (signs) {Разбейте прямую на промежутки;\ проверьте знак внутри каждого};
\node[action,below=of signs] (choose) {Выберите промежутки\ с требуемым знаком};
\node[decision,below=10mm of choose] (strict) {Неравенство строгое?};
\node[side,below left=14mm and 7mm of strict] (yes) {Исключите\ все нули};
\node[side,below right=14mm and 7mm of strict] (no) {Включите\ допустимые нули};
\node[action,below=37mm of strict] (domain) {Исключите недопустимые точки;\ учтите диапазон процесса};
\node[terminal,below=of domain] (end) {Запишите ответ};
\foreach \a/\b in {start/zero,zero/roots,roots/signs,signs/choose,choose/strict,domain/end}
\draw[arrow] (\a) -- (\b);
\draw[arrow] (strict) -- node[above left,font=\small]{да} (yes);
\draw[arrow] (strict) -- node[above right,font=\small]{нет} (no);
\draw[arrow] (yes.south) |- ($(domain.north)+(0,5mm)$) -- (domain.north);
\draw[arrow] (no.south) |- ($(domain.north)+(0,5mm)$) -- (domain.north);
\end{tikzpicture}
\end{center}
\clearpage

\heading{7. Приёмы для дополнительных задач}
\subheading{Куб и четвёртая степень}
\checkitem{Мы учитываем смысл переменной: длина положительна, абсолютная температура положительна. На неотрицательных числах куб и четвёртая степень возрастают. Поэтому сравнение степеней сохраняет порядок чисел.}
Для граничного значения при $x\ge0$ и $A\ge0$:
\[
x^3=A\ \Longrightarrow\ x=\sqrt[3]{A},\qquad
x^4=A\ \Longrightarrow\ x=\sqrt[4]{A}.
\]
\textbf{Примеры.} Решая $x^3=0,064$, получаем $x=0,4$. Поскольку куб возрастает, это максимальное неотрицательное $x$, при котором $x^3$ не больше $0,064$. Решая $x^4=81$ при $x\ge0$, получаем $x=3$. Это минимальное неотрицательное $x$, при котором $x^4$ не меньше $81$.

\checkitem{Мы сначала выражаем степень искомой величины и упрощаем числовой коэффициент, затем извлекаем корень. Например, $5\cdot10^{-8}\cdot4\cdot10^{12}=2\cdot10^5$.}
В формуле натяжения троса физически допустимы $l>0$ и $T\ge0$. Для абсолютной температуры используют кельвины: $500\text{ К}$; знак градуса не ставят.

\subheading{Переменная в знаменателе}
\checkitem{Мы записываем ограничения до умножения на знаменатель. Для приближающегося источника звука в рассматриваемом случае $0\le v<c$, поэтому $1-v/c>0$.}
Для минимальной скорости, при которой частота достигает нужного значения $F>f_0$, решаем граничное уравнение:
\[
\frac{f_0}{1-v/c}=F
\quad\Longrightarrow\quad
v=c\left(1-\frac{f_0}{F}\right).
\]
Здесь $f_0$ и $F$ --- положительные частоты, $c>0$. При увеличении $v$ в диапазоне $0\le v<c$ знаменатель уменьшается, поэтому частота растёт. Равенство даёт минимальную скорость для достижения требуемой частоты. Полученное значение проверяем по ограничению $v<c$.

\textbf{Различие частот.} Для условия «отличаются не менее чем на $\Delta f$» и более высокого второго сигнала граничная частота равна $F=f_0+\Delta f$. Подставляем её в уравнение $f(v)=F$.

\clearpage
\heading{Тренировочный блок — Путь А}
\textbf{Задания 1--5.} Решите самостоятельно. Записывайте основные преобразования и проверяйте единицы. Ответы находятся после всех заданий.
\begin{enumerate}
\item Расстояние, пройденное свободно падающим из состояния покоя телом, задаётся формулой $h(t)=5t^2$, где $t$ --- время в секундах, $h$ --- расстояние в метрах. Найдите $h(0,8)$.
\item Длина стержня при температуре $T$ равна
\[
L(T)=L_0(1+\alpha T),\qquad
L_0=8\text{ м},\quad \alpha=2\cdot10^{-5}\,{}^{\circ}\text{C}^{-1}.
\]
При какой температуре стержень будет на $4$ мм длиннее, чем при $0\,{}^{\circ}\text{C}$?
\item Прибыль от продажи $q$ изделий равна $\Pi(q)=300q-200\,000$ руб., где $q$ --- целое неотрицательное число. Найдите наименьшее число изделий, которое нужно продать для прибыли не менее $300\,000$ руб.
\item Камень падает до поверхности воды; расстояние падения равно $h(t)=5t^2$ м. До дождя падение длилось $0,8$ с, после дождя --- $0,6$ с. На сколько метров поднялся уровень воды?
\item Решите уравнение
\[
\frac{x^2}{12}-\frac{x}{2}=\frac43.
\]
Укажите все действительные корни.
\end{enumerate}
\vfill
\textbf{Проверьте себя.} В задаче 2 переведите миллиметры в метры. В задаче 3 проверьте найденное целое значение и предыдущее. В задаче 4 сравните значения высоты в двух состояниях.

\clearpage
\heading{Тренировочный блок — продолжение}
\textbf{Задания 6--10.} Особое внимание уделите выбору корней и границам применимости формул.
\begin{enumerate}[start=6]
\item Температура прибора после включения изменяется по закону $T(t)=20+20t-t^2$, где $t$ измеряется в минутах, $T$ --- в градусах Цельсия. При первом достижении $84\,{}^{\circ}\text{C}$ прибор автоматически отключается. Через сколько минут это произойдёт?
\item Высота камня над землёй равна $y(x)=-x^2/40+x$, где $x$ --- горизонтальное расстояние в метрах, $0\le x\le40$. Стена имеет высоту $6,5$ м. Найдите наибольшее расстояние $x$, при котором камень находится ровно на $1$ м выше стены.
\item Высота мяча равна $h(t)=10t-t^2$ м при $0\le t\le10$ с. Сколько секунд мяч находится на высоте не менее $21$ м?
\item Для вертикального цилиндрического бака высота воды при сливе равна
\[
h(t)=16\left(1-\frac{t}{400}\right)^2,\qquad 0\le t\le400,
\]
где $h$ --- в метрах, $t$ --- в секундах. Через сколько секунд после открытия крана останется четверть первоначального объёма воды?
\item Температура вещества в течение наблюдения задаётся формулой $T(t)=70+14t-t^2$, где $0\le t\le14$, $t$ --- время в минутах, $T$ --- температура в градусах Цельсия. Укажите все моменты, когда температура не ниже $115\,{}^{\circ}\text{C}$, и найдите продолжительность этого периода.
\end{enumerate}

\clearpage
\heading{Дополнительные задачи}
\textbf{Задания 11--13.} Примените приёмы раздела 7. Проверьте ограничения и найдите требуемый минимум или максимум.
\begin{enumerate}[start=11]
\item \textbf{Подводный зонд.} На верфи инженеры проектируют новый подводный зонд для изучения морских глубин. Конструкция будет крепиться ко дну при помощи троса. Зонд имеет кубическую форму, а значит, сила натяжения троса определяется по формуле:
\[
T=\rho g l^3-mg,
\]
где $l$ --- линейный размер аппарата в метрах, $\rho=1000\text{ кг/м}^3$ --- плотность воды, $g$ --- ускорение свободного падения (считайте $g=10\text{ Н/кг}$), а $m=25\text{ кг}$ --- масса зонда. Каковы могут быть максимальные линейные размеры зонда, чтобы обеспечить его эксплуатацию в условиях, когда сила натяжения троса будет не больше, чем $1000$ Н? Ответ выразите в метрах.

\item \textbf{Температура звезды.} Для определения эффективной температуры звёзд используют закон Стефана--Больцмана, согласно которому мощность излучения нагретого тела прямо пропорциональна площади его поверхности и четвёртой степени температуры:
\[
P=\sigma ST^4,
\]
где $\sigma=5,7\cdot10^{-8}\,\dfrac{\text{Вт}}{\text{м}^2\cdot\text{К}^4}$ --- постоянная, площадь $S$ измеряется в квадратных метрах, температура $T$ --- в кельвинах, а мощность $P$ --- в ваттах. Известно, что некоторая звезда имеет площадь поверхности
\[
S=\frac{1}{228}\cdot10^{20}\text{ м}^2,
\]
а излучаемая ею мощность $P$ не менее $1,5625\cdot10^{25}$ Вт. Определите наименьшую возможную температуру этой звезды. Приведите ответ в кельвинах.
\end{enumerate}

\clearpage
\heading{Дополнительные задачи — продолжение}
\begin{enumerate}[start=13]
\item \textbf{Гудки тепловозов.} Перед отправкой тепловоз издал гудок с частотой $f_0=590$ Гц. Чуть позже издал гудок подъезжающий к платформе тепловоз. Из-за эффекта Доплера частота второго гудка $f$ больше первого: она зависит от скорости тепловоза по закону
\[
f(v)=\frac{f_0}{1-\dfrac vc},
\]
где $c$ --- скорость звука в воздухе (в м/с). Человек, стоящий на платформе, различает сигналы по тону, если они отличаются не менее чем на $10$ Гц. Определите, с какой минимальной скоростью приближался к платформе тепловоз, если человек смог различить сигналы, а $c=300$ м/с. Ответ выразите в м/с.
\end{enumerate}
\vfill
\textbf{Перед сверкой с ответами.} Проверьте единицы и ограничения. Для каждого результата убедитесь, что равенство на границе выполняется и что найденная граница действительно даёт требуемый минимум или максимум.

\clearpage
\heading{Ответы к тренировочному блоку}
\begin{center}
\begin{tabularx}{\linewidth}{@{}cX@{}}
\toprule
№ & Ответ \\
\midrule
1 & $3,2$ м \\
2 & $25\,{}^{\circ}\text{C}$ \\
3 & $1667$ изделий \\
4 & $1,4$ м \\
5 & $x=-2$; $x=8$ \\
6 & $4$ мин \\
7 & $30$ м \\
8 & $4$ с \\
9 & $200$ с \\
10 & $t\in[5;9]$ мин; продолжительность $4$ мин \\
11 & $0,5$ м \\
12 & $5000$ К \\
13 & $5$ м/с \\
\bottomrule
\end{tabularx}
\end{center}
\vspace{7mm}
\textbf{После проверки.} Отметьте в чек-листе освоенные действия. Если ответ отличается, проверьте последовательно: перевод условия, единицы, арифметику, допустимость корней и включение границ.
\end{document}
